\justifying
\NSFCBodyText

\subsubsection{年度研究计划}

研究期限为 2027 年 1 月 1 日至 12 月 31 日。各季度依次完成资料审计、联合模型、约束评估和集成应用，同时尽早启动接口对接与论文组织。

\textbf{1—3 月：建立可检验的数据基础。}
完成 CIESM 变量和频率审计、参考资料处理、地形分层与时间划分，建立插值及统计校正基线。与实验室明确可提供的数据样例、接口和应用方式。阶段交付包括资料清单、样本索引、粗化算子和数据差异报告；若关键频率不足，在本阶段调整数据获取计划并记录核心目标风险。

\textbf{4—6 月：检验联合与时间信息的贡献。}
完成等信息独立模型和联合模型，比较有无地形、单日与相邻日输入，冻结主次评价量及参数选择。形成首轮受控实验、计算预算记录和方法论文初稿。根据降水与风的训练稳定性确定专用输出配置，不以延长训练替代问题诊断。

\textbf{7—9 月：完成约束和外部条件验证。}
完成显式投影、软约束强度分析、地形留出和极端统计；在确定性核心稳定且验证结果支持时，对可选随机残差模型进行概率校准检验。核心任务稳定且资料允许时开展小规模 6 小时条件重建。形成主要实验结果、失效边界报告，并尽早投稿，为审稿和修改预留时间。

\textbf{10—12 月：完成集成应用与成果整理。}
以固定模型处理 CIESM 区域文件，完成公开格式兼容性测试、实验室接口联调和应用示例，按实际接收情况形成应用证据。整理程序、模型配置、数据说明与论文修改材料，提交结题所需的研究总结。论文接收受外部审稿周期影响，将提前投稿并按真实状态报告，不将投稿等同于接收。

\subsubsection{预期研究结果}

拟形成一套 CIESM 多变量时空超分辨率算法、一套可复现的数据处理与训练评估流程，以及新疆典型区域高分辨率气象场试验数据产品。数据产品同时说明参考来源、分辨率含义、覆盖范围、误差特征和限制，区分模型推断与实测资料。

以在所属领域重要学术期刊发表或获得接收的论文不少于 1 篇、成果与 CIESM 集成并实际应用为目标。前者保留发表或接收证明，后者根据接口运行、实际应用和实验室正式接收分别留存证据。算法代码、参数、评价记录和数据获取说明按数据许可与合作约定整理；尚未获许可的资料不列为可公开发布资产。
